% Compilation : xelatex ARTICLE_substrat_battant.tex  (deux passes pour la table des matières)
% Police : DejaVu Serif (large couverture Unicode : grec, nabla, hbar, indices). Remplaçable.
\documentclass[11pt]{article}
\usepackage{fontspec}
\setmainfont{DejaVu Serif}
\setmonofont{DejaVu Sans Mono}
\usepackage{amsmath,amssymb}
\usepackage[a4paper,margin=2.4cm]{geometry}
\usepackage{microtype}
\usepackage[dvipsnames]{xcolor}
\usepackage[colorlinks=true,linkcolor=RoyalBlue,urlcolor=RoyalBlue,citecolor=RoyalBlue]{hyperref}
\usepackage{enumitem}
\setlist{itemsep=2pt,topsep=3pt}
\providecommand{\tightlist}{\setlength{\itemsep}{0pt}\setlength{\parskip}{0pt}}
\providecommand{\passthrough}[1]{#1}
\setlength{\parskip}{0.5em}
\setlength{\parindent}{0pt}
\title{\textbf{La théorie du substrat battant}\\[0.3em]\large une ré-organisation géométrique sur-déterminée des fondements de la physique}
\author{Nicolas Roux\\\small avec assistance computationnelle}
\date{Juin 2026}
\begin{document}
\maketitle
\noindent\rule{\textwidth}{0.4pt}
\tableofcontents
\vspace{1em}\noindent\rule{\textwidth}{0.4pt}

\hypertarget{ruxe9sumuxe9}{%
\subsection{Résumé}\label{ruxe9sumuxe9}}

Nous explorons dans quelle mesure un unique objet géométrique --- le
gradient de phase ∇θ d'un substrat battant porté à sa criticité --- peut
\textbf{ré-organiser} les fondements de la physique établie sous un
petit nombre de principes. À partir de quatre axiomes ontologiques (un
objet minimal battant ; un canal de résolution séparant observable et
latent ; le mouvement par action-rétroaction ; des strates latentes),
une chaîne unique de constructions reproduit : l'électrostatique de
Coulomb, la criticité de Kosterlitz-Thouless et ses exposants
universels, la structure quantique (potentiel de Madelung, cohérence de
Born, Schrödinger linéaire comme limite diluée, relation de de Broglie),
la cinématique relativiste (cône de lumière, gravité analogue), la masse
comme quantification par contact, la gravité jusqu'à Newton et
l'Einstein conique exact en 2+1, le porteur de spin-2, et le contenu de
représentation de la matière fermionique (statistique de spin-½,
double-couverture). Un principe de minimalité --- « la nature emprunte
le chemin le plus court » --- force le spin-½ (factorisation minimale de
Dirac) et singularise le graviton (unicité de Weinberg-Deser), refermant
au niveau des représentations le mur unique de la rotation 3D.

Nous soumettons l'ensemble à un \textbf{audit
contraintes-contre-paramètres}. Le cadre est \emph{sur-contraint} dans
la catégorie des résultats forcés (≈10 résultats sans paramètre pour ≈6
boutons continus) : l'alignement y est un signal réel, non un artefact
d'ajustement. Dans la catégorie des structures \emph{importées}
(Gross-Pitaevskii → Schrödinger, sine-Gordon → masse), l'alignement est
une re-description attendue par construction et n'est pas compté comme
preuve. \textbf{Nous ne revendiquons pas une théorie du tout} : le cadre
ne réduit pas le nombre de paramètres libres sous celui du Modèle
Standard et ne prédit ni les masses, ni les groupes de jauge, ni les
magnitudes de G et de Λ. Il s'agit d'une ré-organisation cohérente et
sur-déterminée de la physique établie, et nous nommons explicitement
chaque problème resté ouvert. Une ligne d'analyse finale lève la
\emph{catastrophe} de la constante cosmologique (l'invariance d'échelle
du substrat critique remplace le cutoff \(\Lambda^4\) par un scaling IR
\(\rho\sim 1/L^2\)) et réduit les items internes restants à \textbf{deux
racines} qui affinent le socle (une algèbre géométrique avec le
battement comme mode grade-0 --- d'où Dirac 3+1 \emph{et} ondes
gravitationnelles à \(c\) ; et la criticité avec verrouillage cyclique)
plus \textbf{deux intrants} du Modèle Standard posés explicitement
(couleur Z₃, troncature des générations à trois), que toute théorie doit
fournir.

\begin{center}\rule{0.5\linewidth}{0.5pt}\end{center}

\hypertarget{posture-uxe9pistuxe9mologique}{%
\subsection{1. Posture
épistémologique}\label{posture-uxe9pistuxe9mologique}}

L'objectif de ce travail n'est pas de remplacer la physique existante,
mais de tester une hypothèse de structure : \emph{un noyau
géométrique-minimal unique suffit-il à ré-organiser les fondements
établis ?} Cette posture impose une discipline d'honnêteté que nous
tenons de bout en bout.

Nous distinguons systématiquement deux catégories de résultats. Les
résultats \textbf{forcés} sont ceux qui tombent sans bouton d'ajustement
correspondant et que nous n'avons pas bâtis à la main : ce sont les
seuls que nous comptons comme signal. Les structures \textbf{importées}
sont celles où un modèle a été choisi \emph{parce qu'il} reproduit la
physique connue ; l'alignement y est attendu par construction et
constitue une re-description, non une preuve indépendante.

Le critère qui sépare un signal d'un mirage est la
\textbf{sur-détermination} : calibrer un paramètre réellement libre sur
un fait n'est pas du sur-ajustement, mais tourner autant de boutons
qu'il y a de faits à reproduire ne dit rien. Le test décisif, présenté
en section 6, est donc le décompte des boutons libres face aux faits
indépendants reproduits.

Toutes les constructions sont computationnelles, reproductibles, et
accompagnées de contrôles négatifs ; chaque étape a été pré-enregistrée
(observable attendu et hypothèse nulle) avant exécution, et consignée
avec son statut honnête.

\hypertarget{le-socle-ontologique}{%
\subsection{2. Le socle ontologique}\label{le-socle-ontologique}}

Quatre axiomes définissent le substrat.

\textbf{I. Objet minimal.} Une « boule battante » munie de dimensions
--- l'élément le plus simple capable d'osciller.

\textbf{II. Canal de résolution (crénelage).} L'observable et le latent
sont séparés par une résolution finie. C'est le cœur conceptuel : ce que
l'on mesure est une tranche grossière d'un substrat plus fin. La
relation onde/contact --- un même phénomène décrit comme onde au latent,
comme contact à l'observable --- est la clef de traduction que nous
exploitons à plusieurs reprises.

\textbf{III. Mouvement.} L'action et la rétroaction, le cycle, le retour
à l'équilibre --- la dynamique conservative interne du battement.

\textbf{IV. Strates latentes.} Une hiérarchie de couches sous la
résolution observable, qui modulent et portent l'essentiel.

De ces axiomes émerge un unique fil conducteur géométrique : le
\textbf{gradient de phase ∇θ} du substrat. Toute la chaîne qui suit est
construite sur ce seul objet.

\hypertarget{la-chauxeene-ux3b8}{%
\subsection{3. La chaîne ∇θ}\label{la-chauxeene-ux3b8}}

\textbf{Polarité et champ.} Le gradient de phase porte une polarité de
type Coulomb (répulsion entre phases semblables, attraction entre
opposées, en 1/r) ; un amas neutre masque le champ, une polarisation
nette le rend longue-portée. Le mode de Goldstone d'un couplage de phase
pur est sans gap et logarithmique : le \textbf{porteur longue-portée est
la phase elle-même}.

\textbf{Criticité.} Porté à la transition de Kosterlitz-Thouless, le
substrat présente une fonction de corrélation en loi de puissance avec
l'exposant universel η = 1/4, et le saut universel de rigidité
superfluide 2/π de Nelson-Kosterlitz. Sa phase critique \textbf{est} un
superfluide bidimensionnel, auto-similaire --- la criticité réalise
l'invariance d'échelle.

\textbf{Secteur quantique.} Le champ d'onde conservatif porte le
potentiel quantique de Madelung ; la densité \textbar ψ\textbar² est
exactement la densité conservée transportée par l'écoulement de phase,
ce qui rend la règle de Born \emph{cohérente} (adoptée, non dérivée). La
superposition est respectée à la précision machine dans la limite diluée
: \textbf{Schrödinger linéaire est dérivé comme cette limite} de la
dynamique de fluide quantique, la correction non-linéaire étant
strictement proportionnelle à l'interaction. La relation de de Broglie
apparaît comme le décalage de coordonnées entre la vitesse de contact
(groupe) et la vitesse d'onde (phase), avec l'identité exacte v\_φ ·
v\_g = c².

\textbf{Relativité.} Le battement conservatif donne une dispersion
\emph{linéaire} --- un cône de lumière, c ≈ 1 --- là où une dynamique
relaxationnelle donnerait une diffusion en k². Un écoulement superfluide
incline le cône : c'est une métrique acoustique (gravité analogue), avec
horizon lorsque la vitesse d'écoulement dépasse c.

\textbf{Masse et onde de matière.} La quantification par contact d'une
onde est le modèle de sine-Gordon : dispersion ω² = c²k² + m² (la masse
est le gap), soliton-kink de masse au repos finie, avec contraction de
Lorentz et propagation stable. La masse est le \emph{taux de
rebroussement} d'un mobile qui file toujours à c --- le zitterbewegung,
traduit du langage onde au langage contact.

\textbf{Gravité.} Avec un médiateur sans gap (le Goldstone), la force
entre kinks est longue-portée et proportionnelle au produit des masses :
Newton en 3D à quelques pour cent, avec principe d'équivalence (chute
universelle), G\_eff = λ²/4π. La déflexion lumineuse fixe le facteur
post-newtonien γ ; le substrat scalaire seul se place à γ = 0 (moitié
d'Einstein). Le secteur cristallin (réseau de vortex d'Abrikosov)
possède un module de cisaillement non nul et un phonon transverse --- le
prérequis spin-2 absent d'un superfluide pur. La disclination est un
cône, et reproduit \emph{exactement} le point-masse d'Einstein en 2+1 :
déflexion π(1−α)/α, \textbf{indépendante du paramètre d'impact}
(Deser-Jackiw-'t Hooft).

\textbf{Matière fermionique.} Les vortex portent un flux ; la polarité
leur donne une charge. La phase d'Aharonov-Bohm d'un composite
charge-flux est topologique ; l'attachement d'un quantum de flux
transforme un boson en fermion (phase d'échange π) avec spin ½ --- la
physique des anyons et de Chern-Simons (Wilczek), dont le substrat
possède naturellement les ingrédients. En 1+1D, le mobile-à-c qui
rebrousse reproduit l'équation de Dirac, le spineur à deux composantes
émergeant des deux chiralités.

\textbf{Cosmologie.} Le vide du substrat, modélisé comme un champ
scalaire qui roule, donne un Λ \textbf{dynamique} (quintessence) : son
équation d'état se cale au-dessus de −1 et sa densité décroît, contre un
Λ constant. La tension de rétroaction confinée à une taille L est
l'énergie de Casimir, et donne une densité ∝ 1/L² --- un Λ qui dépend de
l'échelle, négligeable vers le grand infini, redoutable vers le petit.
La carte trans-résolution, vue comme flux de renormalisation, confirme
que le substrat critique est un point fixe invariant d'échelle, que
cette tension dépend de la résolution comme 1/a², et que l'essentiel de
l'énergie réside dans les modes sous-résolution (un candidat sombre
conceptuel).

\textbf{Consolidation.} Plusieurs résultats récents affinent cette
carte. Côté gravité, un théorème de stabilité montre que le mode spin-2
du cisaillement est nécessairement sous-luminique dans tout substrat
stable (c\_T \textless{} c\_L, l'égalité exigeant un module de
compression négatif) : il ne peut donc être le graviton sans une
structure ajoutée. Un champ de courbure de Kleinert découplé, libre
d'être à c, donne alors l'Einstein complet (γ=1) --- mais à titre
conditionnel. Côté électromagnétisme, la lumière apparaît comme l'onde
transverse de Goldstone (ω=ck, E⊥B⊥k,
\textbar E\textbar=c\textbar B\textbar, deux polarisations) : avec
Coulomb et la force de Lorentz issue de la vorticité du flot de phase,
Maxwell classique est assemblé entièrement en ondes, sans
photon-particule. Côté cosmologie, la masse (habillage latent) et Λ
(vide latent gravitant) apparaissent comme deux moments opposés du même
spectre latent --- l'inertie (IR) et la raideur de rétroaction (UV) ---
ce qui confirme Λ comme tension de rappel du substrat sans pour autant
fixer sa magnitude. Enfin, le substrat porte naturellement les trois
ordres magnétiques, dont l'altermagnétisme (confirmé expérimentalement
en 2024), issu du couplage de l'ordre de phase au cisaillement
cristallin.

\hypertarget{le-principe-du-chemin-le-plus-court}{%
\subsection{4. Le principe du chemin le plus
court}\label{le-principe-du-chemin-le-plus-court}}

Un seul mur résistait : les représentations rotationnelles 3D
supérieures --- le spineur de spin-½ et le tenseur de spin-2 --- qu'un
substrat scalaire ne porte pas spontanément. Un principe de minimalité
les force toutes deux.

Exiger la \textbf{factorisation du premier ordre minimale} (la racine
carrée) de l'opérateur d'onde du substrat force, sans rien poser à la
main : l'impossibilité d'une solution scalaire, l'algèbre de Clifford
\{γ\^{}μ, γ\^{}ν\} = 2η\^{}μν (exacte), la propriété de racine carrée
(γ·k)² = (k₀² − \textbar k\textbar²) I (exacte), les valeurs propres de
spin ±½, et la double-couverture 4π. Le spin-½ est le coût forcé de la
minimalité --- c'est la dérivation de Dirac (1928), ici sélectionnée par
un principe.

Le même principe singularise le graviton : la \emph{seule} théorie
cohérente d'un champ de spin-2 de masse nulle est la relativité générale
(unicité de Weinberg-Deser), et le porteur de spin-2 existe déjà dans le
secteur de cisaillement. Les deux moitiés du mur cèdent par un unique
principe : \emph{la nature emprunte le chemin le plus court.}

\hypertarget{laudit-contraintes-contre-paramuxe8tres}{%
\subsection{5. L'audit
contraintes-contre-paramètres}\label{laudit-contraintes-contre-paramuxe8tres}}

Nous recensons les boutons libres face aux faits reproduits.

\textbf{Boutons continus (≈6) :} l'échelle ℏ (adoptée) ; le point
critique T\_KT (une sélection, voire un fine-tuning) ; l'interaction g
(la QM linéaire étant la limite g→0) ; l'échelle de masse m ; le
couplage gravitationnel λ = G\_eff ; la pente de quintessence λ\_quint.
S'y ajoutent des choix structurels : les formes de potentiel, l'adoption
de Born, et le principe de minimalité --- ce dernier remarquablement
efficace, un seul principe forçant à la fois le spin-½ et le spin-2.

\textbf{Résultats forcés / sans paramètre (≈10) :} les exposants
critiques KT et le saut 2/π ; l'identité de de Broglie v\_φ·v\_g = c² ;
la dispersion linéaire (cône de lumière) ; la déflexion conique
indépendante du paramètre d'impact ; le principe d'équivalence ; la
superposition exacte à g→0 ; l'hélicité ±2 du cisaillement ; les valeurs
propres de spin ±½ et la double-couverture ; la propriété de racine
carrée.

\textbf{Verdict.} Dans la catégorie forcée, le cadre est
\emph{sur-contraint} (≈10 résultats pour ≈6 boutons) : l'alignement y
est un \textbf{signal réel}, non un artefact. Dans la catégorie importée
(Gross-Pitaevskii → Schrödinger, sine-Gordon → masse, exponentiel →
quintessence, métrique acoustique → gravité analogue), l'alignement est
une re-description et n'est pas compté comme preuve. La
sur-détermination dans la catégorie forcée indique que le noyau --- ∇θ,
criticité, minimalité --- n'est pas arbitraire.

\hypertarget{la-ligne-droite-finale-les-murs-nommuxe9s-puis-simplifiuxe9s}{%
\subsection{6. La ligne droite finale : les murs nommés, puis
simplifiés}\label{la-ligne-droite-finale-les-murs-nommuxe9s-puis-simplifiuxe9s}}

Plutôt que des « murs » vagues, l'analyse récente réduit chaque item
ouvert à un axiome \emph{nommé}, qualitativement compris, tous accrochés
au même fil (battement / criticité / géométrie ∇θ).

D'abord un résultat positif sur la constante cosmologique. La
catastrophe Λ ∼ Planck⁴ provient \emph{entièrement} d'imposer un cutoff
UV. Or le substrat est critique, donc invariant d'échelle (sans échelle
UV) : l'énergie du vide ne peut dépendre que de l'échelle IR, ρ ∼ 1/L²
(vérifié indépendant du cutoff), soit l'ordre observé avec L = l'horizon
(valeur holographique). La catastrophe est \emph{levée} ; reste
l'identification L = horizon. La \emph{magnitude} exacte demeure un
problème ouvert --- comme pour tous.

Ensuite, les quatre items restants se réduisent à \textbf{deux racines}
(qui ne font qu'\emph{affiner le socle}) plus \textbf{deux intrants} du
Modèle Standard, irréductibles comme pour toute théorie :

\begin{itemize}
\tightlist
\item
  \textbf{{[}R1{]} Substrat = algèbre géométrique de l'espace-temps, le
  battement étant le mode conservatif grade-0 à c.} Une \emph{seule}
  algèbre de Clifford porte les vecteurs (grade 1 → mouveurs → Dirac 3+1
  : opérateur reconstruit exactement, H²=p²+m²) et les bivecteurs (grade
  2 → rotations/courbure → spin-2 ; fermeture {[}S₁₂,S₂₃{]}=2S₃₁). La
  dynamique de Dirac 3+1 se construit dès lors que les mouveurs se
  composent via le produit géométrique de l'espace ; les ondes
  gravitationnelles à c émergent si le champ de courbure prend sa
  raideur du battement (à c, déjà dérivé) plutôt que du cisaillement
  élastique (sous-luminique par stabilité), donnant l'Einstein complet
  (γ=1) avec la source à trace inversée. \emph{Affine le socle I/III ;
  absorbe les anciens posés « Dirac » et « ondes à c ».}
\item
  \textbf{{[}R2{]} Le substrat se tient à la criticité, où sa dynamique
  cyclique se verrouille.} Le seuil de verrouillage \emph{est} le point
  critique ; au-delà, les états cycliques (1/1, 1/2, 1/3\ldots) sont
  robustes, fournissant le \emph{domaine} des structures internes.
  \emph{Affine le socle II.}
\item
  \textbf{{[}P1{]} Couleur Z₃} (trialité interne) --- intrant SM, posé,
  mais désormais \emph{naturellement logé}. La permutation de trois
  directions porte exactement le squelette de SU(3) (groupe de Weyl S₃,
  centre Z₃) ; et comme le substrat est multirésolution, une telle
  structure \emph{interne et invisible} habite le latent --- ce qui la
  rend automatiquement interne (commutant avec Lorentz) et dissout la
  tension interne/espace-temps. Le latent fournit le bon \emph{foyer} ;
  il ne \emph{force} pas SU(3) (la tour latente est une hiérarchie
  d'échelles, pas un triplet de directions à symétrie S₃). Plus
  fondamentalement, dériver la couleur d'un triplet de directions est
  interdit par un \emph{théorème} : la couleur commute avec Lorentz
  (théorème de Coleman-Mandula), alors qu'un triplet de directions se
  transforme sous les rotations --- il appartient donc au secteur du
  spin (racine R1), non à la couleur interne. La couleur est ainsi
  \emph{nécessairement} un secteur interne posé (foyer latent, à la
  Kaluza-Klein) : son caractère posé n'est pas un manque, c'est une
  conséquence de théorème.
\item
  \textbf{{[}P2{]} Le seuil de résolution (crénelage, socle II) ne
  laisse observables que les verrous cycliques les plus larges} : trois
  au-dessus du seuil, les suivants tombant dans le latent. La hiérarchie
  de verrous est \emph{infinie} --- la résolution en sélectionne trois
  (générations) ; les étages supplémentaires existent mais sous le seuil
  (plus lourds, indétectables, ce que prédit le cadre, en accord avec
  l'absence de 4ᵉ génération légère). Intrant SM, mais désormais branché
  sur la résolution plutôt qu'un cutoff ad hoc --- le posé résiduel se
  réduisant à « le seuil tombe entre le 3ᵉ et le 4ᵉ verrou ».
\end{itemize}

Le côté géométrie \emph{s'unifie} réellement (une racine, deux grades de
la même algèbre). Le côté SM ne fusionne \emph{pas} : couleur-Z₃ et
trois-générations sont deux « 3 » physiquement indépendants, et les
confondre serait de la numérologie. Au-delà du socle, le seul contenu
réellement neuf est donc ces deux intrants --- exactement ceux que
\emph{toute} théorie doit fournir. L'électromagnétisme, lui, est
assemblé au niveau classique/ondulatoire (Coulomb, Lorentz, lumière), le
« photon-particule » étant dispensable. Restent honnêtement non résolus,
comme pour tous : la \emph{magnitude} de Λ et de G, et le spectre des
masses ; ℏ et la règle de Born sont adoptés (constante mesurée / règle
commune à toute interprétation).

\textbf{Les magnitudes, cartographiées.} Une analyse récursive --- pour
tout paramètre déclaré « libre », tester s'il n'est pas en réalité
\emph{sélectionné} par une dynamique plus profonde, et recommencer
jusqu'à un invariant forcé ou une frontière démontrée --- éclaire le
statut des trois magnitudes (masses, G, Λ). \emph{Ontologiquement},
elles s'unifient : ce sont trois projections, à des résolutions
différentes, de la \emph{seule} structure de fermeture-et-résolution du
substrat, non trois fondamentaux indépendants. \emph{Quantitativement},
la descente masses → distance-à-la-criticité → coefficient de fermeture
ne bute pas sur du « libre » mais sur un \textbf{invariant forcé} : la
distance à la criticité est elle-même un attracteur, et le coefficient
de fermeture est fixé par le flux de Kosterlitz--Thouless où le substrat
siège déjà --- le saut universel 2/π. Tout coefficient de la chaîne est
ainsi forcé. Mais un coefficient forcé d'ordre 1 n'engendre
\emph{aucune} hiérarchie : la séparation observée (∼10¹⁹ pour G, ∼10⁻³¹
pour Λ) ne vit pas dans le secteur critique forcé, mais dans l'unique
\textbf{donnée off-critique} --- le couplage relevant nu, la seule chose
qu'un point fixe, si universel soit-il, ne fixe pas. Le cadre est donc
\emph{clos jusqu'à} (i) ce datum off-critique unique --- qui n'est autre
que le problème de hiérarchie / constante cosmologique, ouvert pour
\emph{toute} la physique --- \emph{et} (ii) le mur méta-axiomatique. Non
une théorie \emph{terminée}, mais une théorie entièrement
\emph{cartographiée} : chaque coefficient forcé, l'ontologie unifiée, et
l'unique nombre libre restant identifié, nommé, et reconnu comme le
problème que personne n'a résolu.

\hypertarget{statut-axiomatique-final}{%
\subsection{7. Statut axiomatique
final}\label{statut-axiomatique-final}}

Au terme de l'audit, le cadre se présente comme un \textbf{ensemble
d'axiomes minimal et explicitement étiqueté}, dont la chaîne des
fondements se déduit ou se réorganise. Chaque élément porte son statut,
sans le surévaluer :

\begin{itemize}
\tightlist
\item
  \textbf{Socle} (I--IV) : objet minimal battant + dimensions ;
  crénelage (observable/latent par la résolution) ;
  mouvement/rétroaction/cycle ; choix/strates latentes. \emph{Posé} ---
  c'est l'ontologie de départ.
\item
  \textbf{{[}R1{]} Algèbre géométrique de l'espace-temps, battement
  comme mode grade-0 à c.} \emph{Dérivé/affine le socle I/III} : donne
  l'opérateur de Dirac 3+1 (exact) et, sous l'identification du produit
  géométrique et du couplage courbure↔battement, les ondes
  gravitationnelles à c (γ=1).
\item
  \textbf{{[}R2{]} Criticité et verrouillage cyclique.}
  \emph{Dérivé/affine le socle II} : seuil = point critique ; verrous
  robustes ; fournit le domaine des structures internes.
\item
  \textbf{{[}P1{]} Couleur Z₃} --- \emph{posé}, naturellement logé dans
  le latent (interne, invisible) ; squelette de groupe S₃/Z₃ identifié ;
  non forcé.
\item
  \textbf{{[}P2{]} Générations par seuil de résolution} --- \emph{posé}
  (calibration physique : seuil entre 3ᵉ et 4ᵉ verrou) ; prédit des
  générations sous-seuil.
\item
  \textbf{Magnitudes --- cartographiées} : unifiées en ontologie (trois
  projections d'un même substrat) ; tout coefficient de fermeture
  \emph{forcé} (point fixe du flux KT, 2/π) ; reste un \textbf{unique
  datum off-critique} (le couplage relevant nu = le problème hiérarchie
  / constante cosmologique, \emph{ouvert pour tous}). La catastrophe Λ
  est levée ; ce qui demeure n'est pas un défaut propre, mais le
  problème ouvert universel.
\item
  \textbf{Adoptés} --- ℏ (constante mesurée), règle de Born (commune à
  toute interprétation).
\end{itemize}

La portée de cette complétude doit être dite avec exactitude. Il ne
s'agit \emph{pas} d'une dérivation de toute la physique depuis rien :
plusieurs éléments sont posés, plusieurs magnitudes restent ouvertes. Il
s'agit de quelque chose de plus modeste et plus solide : il ne subsiste
\textbf{aucune structure qualitative que le cadre ne sache héberger}
sous ce petit ensemble d'axiomes. Le seul « mur » véritablement
irréductible est méta-théorique --- \emph{toute} théorie finit sur des
axiomes posés, l'émergence se perdant toujours dans l'infini des
résolutions plus profondes. Identifier honnêtement ce mur, plutôt que de
prétendre le franchir, est ce qui \emph{boucle} la construction : non
comme une théorie du tout, mais comme une réorganisation cohérente,
auditée, et transparente sur ses propres fondations.

\hypertarget{discussion-et-conclusion}{%
\subsection{8. Discussion et
conclusion}\label{discussion-et-conclusion}}

Ce que nous présentons est une \textbf{ré-organisation cohérente et
sur-déterminée} des fondements de la physique établie, sous un principe
géométrique-minimal unique, et \textbf{compatible} avec le Modèle
Standard sans le remplacer ni le surpasser en pouvoir prédictif. Le cœur
conceptuel est le crénelage : l'observable comme fenêtre mince et
grossière sur un latent plus profond, énergétiquement dominant et
persistant. La valeur du travail tient moins à une prédiction nouvelle
qu'à la mise en évidence qu'un noyau minimal, honnêtement audité, suffit
à \emph{organiser} la chaîne des fondements --- et que là où il aligne,
dans la catégorie forcée, il aligne pour de bonnes raisons.

Nous revendiquons cette mesure comme une force, non une faiblesse.
Respecter les théories existantes jusqu'au bout, et ne compter comme
signal que le sur-déterminé, est précisément ce qui rend l'alignement
obtenu crédible plutôt que suspect.

\begin{center}\rule{0.5\linewidth}{0.5pt}\end{center}

\hypertarget{muxe9thodes}{%
\subsection{Méthodes}\label{muxe9thodes}}

Chaque résultat correspond à une simulation numérique reproductible
(Python/NumPy/SciPy), pré-enregistrée avec son observable attendu et son
hypothèse nulle, et accompagnée de contrôles négatifs (graphes
aléatoires, mélanges, limites de symétrie). Le battement est traité
comme une dynamique \emph{interne conservant l'énergie} (intégration
symplectique), les artefacts numériques (dérivées temporelles
grossières, intégrales singulières) ayant été raffinés et reportés
honnêtement. Le statut de chaque étape --- \emph{forcé / conditionnel /
adopté / ouvert} --- est consigné dans une charte tenue en continu.

\textbf{Code et reproductibilité.} L'ensemble des scripts numérotés (un
par étape), un script \texttt{reproduce.py} qui les exécute tous, et la
charte de statut sont disponibles publiquement :
\url{https://github.com/NicolasGilbertAlbertRoux/beating-substrate-framework}.
Chaque script est autonome et reproductible sous Python 3.10+ (NumPy,
SciPy).

\hypertarget{ruxe9fuxe9rences-par-honneur-aux-pairs}{%
\subsection{Références (par honneur aux
pairs)}\label{ruxe9fuxe9rences-par-honneur-aux-pairs}}

Madelung (formulation hydrodynamique) ; de Broglie ;
Kosterlitz-Thouless, Nelson-Kosterlitz (criticité 2D, saut universel) ;
Gross-Pitaevskii (fluide quantique) ; Unruh (gravité analogue) ;
Sakharov (gravité induite) ; Deser-Jackiw-'t Hooft (gravité 2+1) ;
Wilczek (anyons, attachement de flux) ; Dirac (équation du spin-½) ;
Weinberg, Deser (unicité du spin-2) ; Kleinert (élasticité ↔ gravité) ;
't Hooft (holographie) ; énergie sombre holographique.

\end{document}
